% !TEX root = ../../tree_notes.tex
% =====================================================================
%  7.1  The BST Property and Searching          sec:c05-property
%  Exports: def:c05-prop  thm:c05-inorder  prop:c05-search
%           def:c05-bounds  thm:c05-bounds  lst:c05-search  lst:c05-validate
%           ex:c05-prop-run
%           prob:c05-search  lst:c05-search  (practice: LC 653)
% =====================================================================
\section{The BST Property and Searching}\label{sec:c05-property}

\begin{previouslybox}
\begin{enumerate}[nosep, leftmargin=1.6em]
  \item What are the three DFS orders, and which one visits a binary tree's keys
        \emph{left-node-right}? (\chref{c03}, \thmref{c03-one-walk})
  \item Given a \code{TreeNode*} root, how do you descend it with a compare-and-step loop instead
        of recursing into both children? (\chref{c02})
  \item In a sorted array, how many elements can you discard with one comparison? Why is that
        $\bigO{\log n}$? (\chref{c01} intuition; here the array is \emph{shaped} like its own
        comparisons)
  \item What is the inorder sequence of the running BST $B$ of \figref{c05-running}? Read it off
        the picture before you trust the theorem.
\end{enumerate}
\end{previouslybox}

\begin{topicintro}
A binary tree is a shape; a \term{Binary Search Tree} (BST) is a shape \emph{plus a promise about
its contents}. The promise is one line -- $\text{left} < \text{node} < \text{right}$ \emph{at every
node, recursively} -- and it is so strong that it turns the tree into its own search index: a
sorted array that you can \emph{edit}. This section fixes the promise precisely (the recursion is
the part everyone gets wrong), proves the single structural fact that makes BSTs useful --
\thmref{c05-inorder}: the inorder walk \emph{is} the sorted order -- and then shows the one-descent
search that reads that order in $\bigO{h}$ comparisons instead of scanning all $n$ nodes. We close
on the classic trap: validating a BST by checking each node against \emph{only its children}, which
looks right and is wrong, and the $\mathit{lo}/\mathit{hi}$ bound descent that actually works.
\textbf{Prerequisites.} Binary-tree recursion and the \code{TreeNode} struct from Chapter~2
(\chref{c02}); the three traversal orders from Chapter~3 (\chref{c03}), where \emph{inorder} is
defined and never re-derived here. This section adds \emph{the ordering invariant and the descent
it buys you}; it does not re-teach the walk.

\textbf{After this section you can}
\begin{itemize}[nosep]
  \item state the BST invariant as a \emph{recursive} bound, not a local child comparison; \hfill\lv{2}
  \item prove that inorder $\equiv$ sorted in \emph{both} directions and use it as a black box; \hfill\lv{2}--\lv{3}
  \item search, and locate insertion points, with one root-to-leaf descent in $\bigO{h}$; \hfill\lv{2}
  \item validate a candidate BST correctly with a $\mathit{lo}/\mathit{hi}$ bound descent that
        beats the local-check trap. \hfill\lv{3}
\end{itemize}
\end{topicintro}

\begin{recognitionbox}[Pattern card -- ordered content in a binary shape]
\reckey{the statement says ``a \textbf{BST} is given'', or promises ``$\mathit{left} < \mathit{root}
< \mathit{right}$ for every node'', and hands you a \code{TreeNode* root}.}
\recuse{
  \item \emph{does this value exist?} / find where a key belongs -- descend, never scan
        (\probref{c05-search});
  \item ``is this a valid BST?'' -- the bound descent of \thmref{c05-bounds}, \emph{not} a local check;
  \item any ``ordered'' question (kth smallest, closest, min gap) -- \secref{c05-queries} rides the
        sorted line this section proves exists;
  \item building/editing under the rule -- \secref{c05-insert-delete} keeps the invariant alive.
}
\recnotuse{
  \item the contents carry \emph{no order} -- that is a plain binary tree, \chref{c02}, \chref{c03};
  \item you need the height to stay $\bigO{\log n}$ under adversarial inserts -- \chref{c06};
  \item keys are dense integers with many rank queries by index -- Fenwick \chref{c14}.
}
One-line test: \emph{is there a \textbf{sorted order hiding in the structure itself}, so a
comparison tells you which child to follow? Then it is a BST -- descend.}
\end{recognitionbox}

\subsection{Intuition}

Forget ``tree'' for a moment and remember \textbf{binary search}. You stand on a middle element,
compare your target to it, and throw away half. Now imagine the array were not stored end-to-end but
\emph{linked by its own comparisons}: each element points to the half-you-keep as its children. That
is a BST. Standing at the root, one comparison says ``the answer is smaller'' or ``larger,'' and you
descend to exactly the subtree that could contain it -- the other subtree is \emph{provably empty of
the target} and you never look at it.

\figref{c05-running} is the object. To find $7$: compare with $8$, go left (smaller); compare with
$3$, go right (larger); compare with $6$, go right; compare with $7$, found. Four comparisons, four
nodes, and we never touched $1,4,10,13,14$ at all. In a $9$-element \emph{array} we might have
scanned all nine; the tree's shape made the same decisions binary search would make, only the
decisions are \emph{woven into the pointers}.

\begin{figure}[H]
\centering
\begin{tikzpicture}[diagram]
  \node[troot] (n8)  at (0,0)    {$8$};
  \node[tnode] (n3)  at (-3,-1.4) {$3$};
  \node[tnode] (n10) at (3,-1.4)  {$10$};
  \node[tleaf] (n1)  at (-4.2,-2.8) {$1$};
  \node[tnode] (n6)  at (-1.8,-2.8) {$6$};
  \node[tleaf] (n14) at (3,-2.8)  {$14$};
  \node[tleaf] (n4)  at (-2.7,-4.2) {$4$};
  \node[tleaf] (n7)  at (-0.9,-4.2) {$7$};
  \node[tleaf] (n13) at (1.8,-4.2)  {$13$};
  \draw (n8) -- (n3);  \draw (n8) -- (n10);
  \draw (n3) -- (n1);  \draw[tedgehl] (n3) -- (n6);
  \draw[tedgehl] (n6) -- (n4);  \draw[tedgehl] (n6) -- (n7);
  \draw (n10) -- (n14); \draw (n14) -- (n13);
  \draw[tedgehl, line width=1.6pt] (n8) -- (n3);
  \node[annot, align=left, font=\scriptsize] at (5.3,-0.4)
    {search $7$: $8\!\to\!3\!\to\!6\!\to\!7$\\ 4 comparisons\\ touched $5$ nodes\\ \emph{not} $9$};
  \node[annot, align=left, font=\scriptsize] at (-7.0,-2.0)
    {left subtree of $8$\\ $=\{1,3,4,6,7\}$\\ all $<8$\\ \emph{recursively}:\\ node $3$'s right\\ subtree $\{6,7\}$ is\\ $>3$ but still $<8$};
\end{tikzpicture}
\caption{Searching $B$ for $7$ follows one root-to-leaf path (thick edges): $8\to3\to6\to7$, four
comparisons, and it never inspects the five keys $1,4,10,13,14$. The descent is only sound because
the invariant holds \emph{recursively} -- every key under $8$'s left child is below $8$, not merely
below $3$. That recursion is exactly what the naive validation gets wrong.}\label{fig:c05-search}
\end{figure}

The cost is the \emph{path length}, which is at most the height $h$. For $B$, $h=3$ and a search
touches $\le 4$ nodes. This is why ``balanced'' matters so much: a well-shaped BST of $n$ keys has
$h\approx\log_2 n$ and searches in $\bigO{\log n}$, but a badly-shaped one (keys inserted in sorted
order) has $h=n-1$ and degenerates to a linked list -- a warning we feel in \secref{c05-contests}
and that Chapter~8 exists to fix.

\subsection{Definitions and Notation}

\begin{definition}{The BST property (recursive bounds)}{c05-prop}
A binary tree $B$ is a \term{binary search tree} if for every node $v$ with key $k$,
\emph{every} key in $v$'s left subtree is $< k$ and \emph{every} key in $v$'s right subtree is $> k$,
and both subtrees are themselves binary search trees. Equivalently: each node $v$ carries an open
interval $(\mathit{lo},\mathit{hi})$ of legal values with $\mathit{lo}<k<\mathit{hi}$; the left child
inherits $(\mathit{lo},k)$ and the right child inherits $(k,\mathit{hi})$. We take keys
\emph{distinct}; duplicates are discussed in \secref{c05-insert-delete}.
\end{definition}

\noindent\textbf{In plain words.} It is \emph{not} enough that a node beats its left child and loses
to its right child. The whole left \emph{subtree} must lie below the node, and that subtree's own
children must respect \emph{both} their local parent and every bound inherited from the ancestors
above. The interval language is the one you will actually test against -- it is what turns
validation into a straightforward descent.

\begin{theorem}{Inorder is the sorted order (both directions)}{c05-inorder}
For a binary tree $B$ with distinct keys, $B$ is a BST \emph{if and only if} its inorder traversal
(left, node, right) lists the keys in strictly increasing order.
\end{theorem}
\begin{proofbox}
$(\Rightarrow)$ Induct on the tree. For a leaf the inorder list is one key -- sorted. Otherwise the
root has key $k$ with left subtree $L$ (all keys $<k$ by \defref{c05-prop}) and right subtree $R$
(all keys $>k$). By induction $\mathrm{in}(L)$ is sorted and $<k$, and $\mathrm{in}(R)$ is sorted and
$>k$, so $\mathrm{in}(L),k,\mathrm{in}(R)$ is one increasing run. $(\Leftarrow)$ Contrapositive: if
$B$ is \emph{not} a BST, some node $v$ with key $k$ has a left-subtree key $x>k$ (or a right-subtree
key $y<k$). Every key of $v$'s left subtree is visited \emph{before} $k$ in inorder, so $x$ appears
before $k$ with $x>k$ -- the inorder list is not increasing. Hence a non-sorted inorder witnesses a
non-BST, proving the converse. \end{proofbox}
\begin{insightbox}
Read the theorem as a \emph{two-way tool}. Forward: to answer any ordered question, run inorder and
you hold a sorted array -- no re-sorting, $\bigO{n}$. Backward: to \emph{check} a tree is a BST, you
may verify its inorder is increasing -- but do it streaming (compare each key to its predecessor)
so you stop at the first descent and never store the whole list. \thmref{c05-bounds} gives an even
cheaper $\bigO{n}$ check that needs no list at all.
\end{insightbox}

\begin{proposition}{One-descent search}{c05-search}
Searching a BST $B$ of height $h$ for a key $x$ visits at most $h+1$ nodes and uses at most $h$
comparisons: start at the root; if $x$ equals the current key stop; if $x$ is smaller descend left;
if larger descend right; a null link means $x\notin B$. The procedure is correct and does not inspect
any node off the root-to-node path.
\end{proposition}
\begin{proofbox}
Each step performs one comparison and moves strictly down, so it halts in $\le h+1$ nodes. By
\defref{c05-prop}, when $x<k$ no key in the right subtree can equal $x$ (they are all $>k>x$), so
discarding it is safe; symmetrically for $x>k$. Thus the target, if present, always stays in the
subtree we descend into, and we stop exactly when we reach it or fall off the tree. \end{proofbox}

\begin{theorem}{Correct validation by bound descent (the local-check trap)}{c05-bounds}
A tree is a BST if and only if every node $v$ with key $k$ satisfies $\mathit{lo}<k<\mathit{hi}$ for
the interval $(\mathit{lo},\mathit{hi})$ inherited from its ancestors, propagated as
$(\mathit{lo},k)$ to the left child and $(k,\mathit{hi})$ to the right child. Checking
\emph{only} $v$ against its immediate children is \textbf{not} sufficient.
\end{theorem}
\begin{proofbox}
The bound descent checks exactly \defref{c05-prop}: an inherited $\mathit{lo}$ is the largest key
that is a proper ancestor with $v$ in its right subtree, and $\mathit{hi}$ the smallest ancestor with
$v$ in its left; $\mathit{lo}<k<\mathit{hi}$ forces $k$ above every left-side requirement and below
every right-side one, i.e.\ the whole subtree respects all ancestors. \emph{Local is not enough:}
in \figref{c05-trap} the node $1$ is $<7$ and in $7$'s left subtree, and $7$ is $>5$ in $5$'s right
subtree -- every node beats its own children, yet $1$ sits in the right subtree of $5$ with $1<5$,
violating \defref{c05-prop}. The bound descent catches it: node $1$ inherits $(5,7)$ and $1\not>5$.
\end{proofbox}

\begin{figure}[H]
\centering
\begin{tikzpicture}[diagram]
  \node[boxnode, draw=warningred, text=warningred] (a) at (0,0) {$5$};
  \node[tnode] (l) at (-2.2,-1.4) {$2$};
  \node[tnode] (r) at (2.2,-1.4) {$7$};
  \node[tleaf, draw=warningred] (x) at (1.1,-2.8) {$1$};
  \node[tleaf] (y) at (3.3,-2.8) {$9$};
  \draw (a) -- (l); \draw (a) -- (r);
  \draw[warningred, line width=1.4pt] (r) -- (x);
  \draw (r) -- (y);
  \node[annot, align=left, font=\scriptsize] at (6.0,-1.6)
    {local check \emph{passes}:\\ $2<5<7$ \checkmark,\\ $1<7<9$ \checkmark\\[-2pt]
     global fails:\\ $1$ is under $5$'s\\ right subtree but $1<5$\\ ✗ \defref{c05-prop}};
\end{tikzpicture}
\caption{The local-check trap. Every node is correctly ordered against \emph{its own children}, yet
the tree is \textbf{not} a BST: node $1$ lies in the right subtree of $5$ (so must exceed $5$) but
equals $1$. Only the inherited-bound test of \thmref{c05-bounds} -- giving node $1$ the interval
$(5,7)$ -- flags it. \probref{c05-search} uses this same $B$ for search; \thmref{c05-bounds} is the
reason validation is a \emph{descent}, not a per-node comparison.}\label{fig:c05-trap}
\end{figure}

\subsection{The algorithm: descend once, carrying the bounds}

Two loops, one shape. Search compares and steps toward one child (\coderef{c05-search});
validation walks the same pointers but carries the $(\mathit{lo},\mathit{hi})$ interval so a deep
node can be checked against \emph{every} ancestor, not just its parent (\coderef{c05-validate}).
Both cost $\bigO{h}$ per operation on the search and $\bigO{n}$ for the full validation (each node
once).

\begin{lstlisting}[caption={Search a BST for a value -- iterative one-descent (LeetCode 700)}, label={lst:c05-search}]
#include <bits/stdc++.h>
using namespace std;

struct TreeNode { int val; TreeNode *left, *right;
                  TreeNode(int v) : val(v), left(nullptr), right(nullptr) {} };

// Returns the node whose val == target, or nullptr. One root-to-leaf path.
TreeNode* searchBST(TreeNode* root, int target) {
    TreeNode* cur = root;
    while (cur != nullptr && cur->val != target) {
        cur = (target < cur->val) ? cur->left   // target smaller -> left subtree
                                  : cur->right; // target larger  -> right subtree
    }
    return cur;   // nullptr if the key is absent
}
\end{lstlisting}

\begin{lstlisting}[caption={Validate a candidate BST with the bound descent -- the fix for the local-check trap}, label={lst:c05-validate}]
#include <climits>

// Every key under `node` must lie strictly inside (lo, hi).
// Use long to survive INT_MIN/INT_MAX keys at the root's own bounds.
bool validRange(TreeNode* node, long lo, long hi) {
    if (node == nullptr) return true;
    if (node->val <= lo || node->val >= hi) return false;   // violates an inherited bound
    return validRange(node->left,  lo,      node->val)      // left must be < node->val
        && validRange(node->right, node->val, hi);          // right must be > node->val
}

bool isValidBST(TreeNode* root) {
    return validRange(root, LONG_MIN, LONG_MAX);
}
\end{lstlisting}
\begin{complexitybox}
Search: $\bigO{h}$ time, $\bigO{1}$ extra space (the loop keeps only one pointer). Validation:
$\bigO{n}$ time -- each node is checked once -- and $\bigO{h}$ stack space for the recursion (worst
$\bigO{n}$ when the tree is a chain). Verdict: fast when $h$ is small, and \emph{nothing} here
assumes balance -- the cost is honestly written in $h$, which is the chapter's running caution.
\end{complexitybox}

\subsection{Worked Example}

\begin{example}{Searching $B$ for present and absent keys}{c05-prop-run}
Take the running BST $B$ (\figref{c05-running}) with inorder $1,3,4,6,7,8,10,13,14$.
\emph{Search for $7$:} start at $8$; $7<8$ go left to $3$; $7>3$ go right to $6$; $7>6$ go right to
$7$; $7=7$ stop. Path $8\to3\to6\to7$, four nodes, three ``which way'' comparisons, and one match.
\emph{Search for $5$ (absent):} $8$; $5<8$ left to $3$; $5>3$ right to $6$; $5<6$ left to $4$;
$5>4$ right $\to$ null. Path $8\to3\to6\to4$, then a null link $\Rightarrow$ report ``absent''. Both
searches never visited $10,13,14$ (nor $1$) -- \propref{c05-search} guarantees the target could not
hide there, and indeed the sorted line shows $5$ sits between $4$ and $6$, which we walked past.
\end{example}
\begin{tracebox}[The bound descent flags node $1$ in the trap tree \figref{c05-trap}]
Validate $5\!\to\!(2,7)$, $7\!\to\!(1,9)$. Open at $5$ with $(\text{-}\infty,+\infty)$: ok.
Descend left to $2$ with $(-\infty,5)$: $2$ in range, its children null, ok. Descend right to $7$
with $(5,+\infty)$: $7>5$ ok. Descend to $7$'s left child $1$ with $(5,7)$: check $1>5$? \textbf{No}
-- return false immediately. The local check would have compared $1$ only to $7$ ($1<7$, ``fine'')
and missed the violation; the inherited lower bound $5$ is what \defref{c05-prop} demands and what
\thmref{c05-bounds} enforces. Tree rejected at the first bad node, no full list built.
\end{tracebox}

\subsection{Variations and Twists}

\begin{itemize}[nosep]
  \item \textbf{Recursive search}: \code{if (!root||root->val==target) return root;} then tail-call
        into one child. Same $\bigO{h}$, but now $\bigO{h}$ stack. \hfill\lv{2}
  \item \textbf{Insertion point}: run the search loop until \code{cur} is null; the last non-null
        node you compared is where the new key attaches. This is \emph{exactly} \defref{c05-insert}
        of \secref{c05-insert-delete} -- search half-written. \hfill\lv{2}
  \item \textbf{Range count $\mathbf{[L,R]}$}: descend, and at a node fully inside the range add its
        subtree size; when a node is below $L$ go right only; above $R$ go left only. Turns a scan
        into $\bigO{h + k}$ -- developed in \secref{c05-queries}. \hfill\lv{3}
  \item \textbf{Streaming validation}: instead of intervals, do an inorder walk and assert each key
        exceeds its predecessor. Equivalent by \thmref{c05-inorder}; needs $\bigO{h}$ and one
        \code{prev} variable. \hfill\lv{3}
  \item \textbf{Duplicates}: allow $x\ge k$ on the left and stop ``found'' at equality; pick a
        convention (all duplicates one side) and keep \thmref{c05-bounds}'s strict/loose signs
        consistent with it. \hfill\lv{3}
\end{itemize}

\subsection{Common Pitfalls}\label{sec:c05-property-pitfalls}

\begin{warningbox}[Trap 1 -- validating with local child comparisons]
The single most common BST bug: assert $\mathit{left->val}<v<\mathit{right->val}$ at each node and
declare victory. \figref{c05-trap} passes that test yet is not a BST. Carry the inherited interval
(\thmref{c05-bounds}, \coderef{c05-validate}) or compare to the inorder predecessor -- never compare
a node only to its children.
\end{warningbox}
\begin{lstlisting}[caption={\textcolor{warningred}{WRONG}: local check -- passes the trap tree of \figref{c05-trap}}, language=C++, basicstyle=\ttfamily\small]
bool ok(TreeNode* r) {
    if (!r) return true;
    if (r->left  && r->left->val  >= r->val) return false;   // only the child...
    if (r->right && r->right->val <= r->val) return false;   // ...not the subtree
    return ok(r->left) && ok(r->right);                       // MISSES ancestor bounds
}
\end{lstlisting}
\begin{lstlisting}[caption={\textcolor{successgreen}{RIGHT}: check every node against inherited bounds}, language=C++, basicstyle=\ttfamily\small]
bool ok(TreeNode* r, long lo, long hi) {
    if (!r) return true;
    if (r->val <= lo || r->val >= hi) return false;      // violates ANY ancestor bound
    return ok(r->left, lo, r->val) && ok(r->right, r->val, hi);
}
\end{lstlisting}

\begin{warningbox}[Trap 2 -- overflow at the sentinel bounds]
Seeding the root's interval with \code{INT\_MIN}/\code{INT\_MAX} and using \emph{strict} inequalities
rejects a legitimate BST whose root actually \emph{equals} \code{INT\_MIN} (or \code{INT\_MAX}),
because the node is tested against its own sentinel. Widen the bounds to \code{long} (as in
\coderef{c05-validate}) or pass $\mathit{lo}/\mathit{hi}$ as \code{nullable} flags instead of extreme
values.
\end{warningbox}
\begin{lstlisting}[caption={\textcolor{warningred}{WRONG}: extreme \code{int} bounds reject a valid single-node tree \{INT\_MIN\}}, language=C++, basicstyle=\ttfamily\small]
return validRange(root, INT_MIN, INT_MAX);   // node == INT_MIN fails "lo < val"
\end{lstlisting}
\begin{lstlisting}[caption={\textcolor{successgreen}{RIGHT}: long bounds strictly outside every int key}, language=C++, basicstyle=\ttfamily\small]
return validRange(root, (long)INT_MIN - 1, (long)INT_MAX + 1);
\end{lstlisting}

\begin{warningbox}[Trap 3 -- scanning the whole tree when the order was free]
Once \thmref{c05-inorder} tells you the structure \emph{is} sorted, walking both children to answer
``is $x$ present?'' throws the invariant away: you pay $\bigO{n}$ for what \propref{c05-search}
does in $\bigO{h}$. If your routine recurses into \emph{both} subtrees on a membership/insertion
question, stop -- one comparison decides which child could hold $x$.
\end{warningbox}

\subsection{Problem Lab}

The section's anchor is LC 700 -- the purest possible BST task: given a BST and a value, return the
subtree rooted at the matching node, or nothing. There is no edit, no construction, no ordering to
prove: it is a single graded test that you used \emph{structure} (one descent) rather than
\emph{brute force} (a full traversal). Everything later in the chapter is this walk with extra
responsibilities.

% ---------------------------------------------------------------------------
\solvedproblem{c05-search}{Search in a Binary Search Tree}{LeetCode 700}
  {https://leetcode.com/problems/search-in-a-binary-search-tree/}{LC Easy L2}
  {BST invariant, one-descent search, iterative vs recursive descent}

\begin{problemstatement}[Search in a Binary Search Tree]
You are given the \code{root} of a BST and an integer \code{val}. Return the node whose
\code{val} equals the given value and the subtree rooted there (as a \code{TreeNode*}), or
\code{nullptr} if no such node exists. The input is \textbf{guaranteed to be a valid BST}. There are
$1\le n\le 2000$ nodes and $1\le \text{val},\ \text{node.val}\le 10^{4}$, all distinct.
\end{problemstatement}

\begin{attemptgate}[Attempt it yourself first -- \lv{2}, give yourself 10 minutes]
This is the chapter's hello-world: it tests only whether you trust the invariant.
\peek{You are told it \emph{is} a BST. That means at any node, \code{val} can live in \emph{at most
one} child. Which one, and how do you decide with a single comparison?}
\nudge{Start at \code{root}. If \code{root->val == val} you are done. If \code{val < root->val} the
answer (if any) is in the left subtree; otherwise the right. Repeat on that child. A null child means
absent -- return \code{nullptr}.}
\fullhint{Set \code{cur = root}. Loop while \code{cur} is non-null and \code{cur->val != val},\newline
\hspace*{1em}stepping \code{cur} to \code{cur->left} if \code{val < cur->val}, else \code{cur->right}.
Return \code{cur}. That is \coderef{c05-search} verbatim -- $\bigO{h}$, no recursion,
no scanning. Note you must return the \emph{node}, not a boolean: the caller wants the whole subtree.}
\end{attemptgate}

\solstep{Understanding the Problem}
We are handed a valid BST -- the invariant of \defref{c05-prop} is a \emph{gift}, not something to
verify. Our job is a membership lookup that hands back the located node so its subtree can be reused.
The trap in the two-star framing is over-thinking: a full DFS that compares \code{val} at every one
of the $n$ nodes is \emph{correct} but ignores the gift and costs $\bigO{n}$. The intended answer is
the $\bigO{h}$ descent of \propref{c05-search}. Return \code{nullptr} when we fall off the tree.

\solstep{Key Observations}
\begin{observation}
Because the tree is a valid BST, at every node $v$ the value \code{val} can appear in \emph{at most
one} subtree: left if \code{val < v->val}, right if \code{val > v->val}, and nowhere if we reach a
null link. So one comparison \emph{eliminates} a whole subtree -- exactly binary search.
\end{observation}
\begin{observation}
The problem wants the \emph{node}, not a yes/no. Stop the descent at the match and return the
pointer; that pointer \emph{is} the rooted subtree requested.
\end{observation}
\begin{observation}
No auxiliary storage is needed: a single moving pointer suffices, so the iterative form is
$\bigO{1}$ extra space and immune to stack overflow on a deep (unbalanced) BST.
\end{observation}

\solstep{Think It Through With Me}
\begin{thinkbox}
\deadend{Traverse the whole tree and return the first node equal to \code{val}.}
{Correct, but $\bigO{n}$: you are treating the data as an unordered binary tree and ignoring the one
promise the statement made. On a chain of $2000$ nodes you would walk $2000$ comparisons where a
smart descent also walks $\sim n$ -- but on a balanced tree you'd walk $n$ where $\bigO{\log n}$
would do. The problem is testing that you \emph{notice} it is a BST.}
\deadend{Recurse into whichever child is indicated, but also ``just in case'' search the other.}
{If you ever search both children, the invariant bought you nothing -- that is the $\bigO{n}$ scan
again, dressed as recursion. One comparison, one child. \propref{c05-search} is a \emph{single}
path.}
\aha{Walk down one path. Compare, pick the only child that could hold \code{val}, move there, and
either hit the value (return the node) or hit null (return \code{nullptr}). $\bigO{h}$, $\bigO{1}$
space, no scanning.}
\end{thinkbox}

\solstep{Solution}
Set \code{cur = root}. Loop while \code{cur} is non-null and \code{cur->val != val}: each iteration
replaces \code{cur} by \code{cur->left} when \code{val < cur->val}, else \code{cur->right}. When the
loop exits, \code{cur} is either the matching node or null; return it. Nothing else is touched, which
is precisely the guarantee of \propref{c05-search}: the target, if present, was always inside the
subtree we descended into.

\solstep{Proof of Correctness}
\begin{proofbox}
We show the loop returns the node with \code{val} if it exists and \code{nullptr} otherwise.
\emph{Invariant:} at the top of each iteration, if a node with value \code{val} exists in the tree, it
lies in the subtree rooted at \code{cur}. Initially \code{cur} is the root -- trivially true. Each
step: if \code{cur->val == val} we exit and \code{cur} is the answer. Otherwise, by \defref{c05-prop},
when \code{val < cur->val} all keys in \code{cur}'s right subtree exceed \code{cur->val > val}, so no
right-subtree node equals \code{val}; the target, if anywhere, is in the left subtree, and setting
\code{cur = cur->left} preserves the invariant. The case \code{val > cur->val} is symmetric.
\emph{Termination:} \code{cur} strictly descends, so after at most $h+1$ steps it is null or a match.
If it becomes null, the invariant forces that no matching node ever existed, so returning null is
correct; if it equals \code{val}, that node is the (unique, keys distinct) match. Hence the returned
pointer is exactly the required subtree root, or null. \end{proofbox}

\solstep{Implementation}
\begin{lstlisting}[caption={Search in a BST -- the graded, submittable solution (LeetCode 700)}, label={lst:c05-search-soln}]
/**
 * Definition for a binary tree node.
 * struct TreeNode { int val; TreeNode *left, *right; };
 */
class Solution {
public:
    TreeNode* searchBST(TreeNode* root, int val) {
        TreeNode* cur = root;
        while (cur != nullptr && cur->val != val)
            cur = (val < cur->val) ? cur->left : cur->right;  // one comparison, one child
        return cur;   // matching node, or nullptr if absent -- O(h), never O(n)
    }
};
\end{lstlisting}

\solstep{Trace}
\begin{tracebox}[Running BST $B$ (\figref{c05-running}) : search 7 (found) and 5 (absent)]
\emph{val=7:} \code{cur}$=8$, $7<8$ so \code{cur}$=3$; $7>3$ so \code{cur}$=6$; $7>6$ so
\code{cur}$=7$; \code{cur->val==7} exit $\Rightarrow$ return the node $7$ (its subtree is just itself).
\checkmark\\
\emph{val=5:} \code{cur}$=8\to3\to6\to4$ (choices $L,R,L,R$); \code{cur->right} is null so \code{cur}
becomes null, loop exits, return \code{nullptr}. \checkmark Keys $10,13,14,1$ were never visited in
either search -- the descent's whole point.
\end{tracebox}

\solstep{Complexity}
\begin{complexitybox}
Time $\bigO{h}$, where $h$ is the height ($\le n$): one comparison per node on a single root-to-leaf
path. Space $\bigO{1}$ for the iterative pointer. For a \emph{balanced} BST $h=\bigO{\log n}$ so this
is binary search; for the worst-case chain (keys inserted in sorted order) $h=\bigO{n}$ and it
degrades to a linear scan -- the very cliffhanger \secref{c05-contests} raises and \chref{c06}
resolves. Verdict: \ac.
\end{complexitybox}

\solstep{Pitfalls and Variants}
\begin{warningbox}
Three ways to lose marks on a two-star problem: (i) writing a $\bigO{n}$ DFS that ``works'' but
ignores the BST promise -- the grader's hidden test on a large balanced tree is a timing tell.
(ii) returning a \code{bool} instead of the node (the caller needs the subtree). (iii) dereferencing
\code{cur->left} before checking \code{cur} is null. \emph{Variants:} return the \code{parent} at the
insertion point and you have \defref{c05-insert} (\secref{c05-insert-delete}); return the last node
$\le$ \code{val} and you have \term{floor}, the seed of the successor work in \secref{c05-queries}.
\end{warningbox}

\begin{insightbox}[Key takeaways]
\begin{enumerate}[nosep, leftmargin=1.6em]
  \item \textbf{The invariant is a licence to discard.} One comparison removes a whole subtree --
        that is \propref{c05-search}, and it is the difference between $\bigO{h}$ and $\bigO{n}$.
  \item \textbf{Trust ``it is a BST''} -- do not validate it here (that is \thmref{c05-bounds}, a
        different task). Use it.
  \item \textbf{Cost is written in $h$, not $n$.} Every chapter-7 routine is $\bigO{h}$; keeping $h$
        small is Chapter~8's job, and the Summit (\probref{c05-summit}) makes it the enemy.
\end{enumerate}
\end{insightbox}

\begin{hardeningbox}
\begin{itemize}[nosep, leftmargin=1.4em]
  \item \textbf{Recursive form}: one tail call into the chosen child; identical complexity but
        $\bigO{h}$ stack -- the iterative loop is the robust default on deep trees.
  \item \textbf{Closest value}: track the best node seen while descending toward \code{val}; when
        you would recurse into a subtree that cannot improve, the running min already covers it --
        the technique of LC 270 (\term{Closest BST Value}), an unlisted cousin.
  \item \textbf{Insertion point}: the loop's final \code{cur} is null, but the \emph{previous}
        pointer is where a new key hangs -- reuse the exact code for \secref{c05-insert-delete}.
\end{itemize}
\end{hardeningbox}

\subsection{Practice Problems}

\begin{stuckbox}
\begin{itemize}[nosep, leftmargin=1.4em]
  \item Did you \emph{descend one path}, not traverse both children? (\propref{c05-search}.)
  \item Did you return the \emph{node} (subtree), not a boolean? (\probref{c05-search}.)
  \item For the practice: are you exploiting sortedness, or secretly $\bigO{n}$ scanning a list of
        values? The BST property should prune, not just decorate.
  \item If validating, did you carry $(\mathit{lo},\mathit{hi})$ rather than compare to children?
        (\thmref{c05-bounds}, \coderef{c05-validate}.)
\end{itemize}
\end{stuckbox}

\begin{practiceset}
  \practice{Two Sum IV -- Input is a BST}{LeetCode 653}
    {https://leetcode.com/problems/two-sum-iv-input-is-a-bst/}{LC Easy L2}
    {Given a BST and a \code{target}, decide whether two \emph{distinct} nodes sum to it. The BST
    promise says the inorder list is sorted (\thmref{c05-inorder}) -- so you already know the
    answer for a sorted \emph{array}: two pointers, or, per value $x$, search for $\mathit{target}-x$
    with \propref{c05-search}. Do \emph{not} dump the tree into a vector and forget the order you
    were gifted; either walk inorder and use a sorted-array two-pointer, or for each node run a
    second $\bigO{h}$ search for its complement. The trap is double-counting the same node when
    $2x=\mathit{target}$ -- require distinct nodes.}
\end{practiceset}

\begin{recallbox}[10-second refresher -- 7.1 The BST Property and Searching]
A BST obeys $\text{left}<\text{node}<\text{right}$ \emph{recursively} (\defref{c05-prop}), so its
inorder walk \emph{is} the sorted order and vice versa (\thmref{c05-inorder}). Search/insert-point
finders descend \emph{one path} in $\bigO{h}$ (\propref{c05-search}, \coderef{c05-search}). Validate
with an inherited $(\mathit{lo},\mathit{hi})$ interval, never a local child comparison
(\thmref{c05-bounds}, \coderef{c05-validate}) -- the trap tree of \figref{c05-trap} looks locally fine
but is not a BST. Cost is written in $h$; keeping $h$ small is \chref{c06}.
\selftest{
  \item State the BST property for the \emph{whole left subtree} of a node, not just its left child.
  \item Prove that a tree whose inorder is strictly increasing must satisfy \defref{c05-prop}.
  \item On $B$ (\figref{c05-running}), trace the search for $10$ and say which nodes are \emph{not}
        visited.
  \item Give a two-node-per-side tree that passes the local check but fails \thmref{c05-bounds}.
  \item Why is seeding validation bounds with \code{int} extremes a bug, and what is the fix?
}
\end{recallbox}
